\begin{comment}
\subsection{Introduction}

A substantial portion of the gamma-ray radiation that we observe from Earth is produced inside our own galaxy, mostly along the galactic plane. At high Galactic latitudes, however, the gamma-ray emission is mostly of cosmological origin.
This extragalactic background radiation (EGB) is the sum of the emission from all the extragalactic gamma-ray sources \cite{fornasa:2015,Fermi-LAT:2014ryh}. Most of the EGB is originated by various types of astrophysical sources which are seen from our viewpoint as point sources. A relevant observable is therefore their differential source-count distribution $dN/dS$ which counts the number of sources at a given integral source flux $S$ (for gamma-ray energies in a specific interval). For bright sources, which are observationally identified and therefore cataloged, the source-count distribution is directly measured by counting the objects in the catalog,
at least above the threshold for which the catalog detection efficiency is equal to 100\%.  
Below this flux threshold, where the detection efficiency is less than 100\%, dedicated Monte Carlo simulations are required to accurately determine the efficiency and use it as a correction to reconstruct the true underlying $\dnds$ \cite{Fermi-LAT:2010tsy,Fermi-LAT:2015otn,Ajello:2015mfa,Marcotulli:2020fpm}.
For all those sources which are too faint to be resolved, their cumulative distribution of photons in the sky defines an almost isotropic cosmic field, conventionally called the isotropic diffuse gamma-ray background (IGRB) or, more precisely, the unresolved gamma-ray background (UGRB). Even though individual sources below the detector flux-threshold cannot be individually seen, it has been shown \cite{Malyshev_2011,Cuoco-1pdf,Zechlin:2016pme,Lisanti:2016jub}  nevertheless, that it is possible to infer their source-count distribution even in this regime, looking at the collective effects of these unresolved sources. This technique, called pixel-count distribution (or 1-point PDF) and pioneered for gamma-rays in \cite{Malyshev_2011}, has been improved in \cite{Cuoco-1pdf,Zechlin:2016pme,Lisanti:2016jub} by employing a pixel-dependent approach, in order to fully explore all the available information and to incorporate the morphological variation of the gamma-ray emission components. Ref. \cite{Cuoco-1pdf} used the first 6 years of \fermi \cite{Atwood:2009ez} data to measure the $dN/dS$ for photons in the energy range (1,10) GeV down to fluxes about one order of magnitude below the \fermi detection threshold. In \cite{Zechlin:2016pme,Lisanti:2016jub} the same technique was used to extend the measurement of the $dN/dS$ to several energy bands between 1 and 171 GeV, thus providing information on the energy dependence of the source-count distribution.
This technique has then been used to study the contribution of individual source classes of emitters, in particular blazars \cite{Korsmeier:2022cwp,Manconi:2019ynl,DiMauro:2017ing,Zechlin:2017uzo}.
The same methodology has also been used to characterize the properties of the unresolved sources in the Galactic Center region in relation to a possible signal from dark matter annihilation \cite{Mishra-Sharma:2016gis,Bariuan:2022phe,Lee:2014mza,Lee:2015fea,Linden:2016rcf,Leane:2019xiy,Chang:2019ars,Buschmann:2020adf,Leane:2020nmi,Leane:2020pfc}.
{  A further methodology called Compound Poisson Generator has also been developed \cite{Collin:2021ufc}, in order to handle biases possibly present in the 1-point PDF method.}
\end{comment}


In this paper we update the measurement of the $dN/dS$ below the detection threshold to the increased statistics offered by 14 years of \fermi data. However, differently from Refs. \cite{Cuoco-1pdf,Zechlin:2016pme}, we adopt here a method based on machine learning techniques to obtain the $dN/dS$ below the detection threshold. 
A similar methodology has also been used recently to investigate gamma-ray unresolved sources close to the Galactic center \cite{Mishra-Sharma:2021oxe,List:2021aer,List:2020mzd}.  
We train a convolutional neural network (CNN) on synthetic gamma-ray maps, built from a wide variety of source-count distributions and then apply the trained CNN to the  14-year \fermi map for photon energies in the (1,10) GeV band. We show that the CNN is able to reconstruct the $dN/dS$, thus obtaining an updated result which is fully compatible with the one obtained in \cite{Cuoco-1pdf}. The methodology presented here is also meant to be a proof of principle for the adoption of a CNN to the reconstruction of the source-count distribution of the extragalactic sky, with the future aim of properly investigating additional features of the $dN/dS$, like energy correlations and the presence of additional components which could be traced to dark matter.
{  An advantage of the CNN method is that it
avoids the need to calculate complicated and numerically demanding likelihoods, as in the 1-point PDF method.
Furthermore, we will describe an improved version of the treatment of a CNN on a spherical domain which will further optimize the computational aspect.}


Data selection and \fermi map generation is discussed in Section \ref{sec:dataselection}, while Section \ref{sec:generation} discusses in detail how synthetic maps for the CNN training and validation are constructed. Section \ref{sec:architechture-and-training} discusses the neural network architecture that we use and its implementation, including validation and error estimation. Section \ref{sec:dnds:results} discusses the analysis of the \fermi map with the trained CNN and presents the ensuing results for the $dN/dS$. Section \ref{sec:dnds:conclusions} gives our conclusions.




%-----------------------

